\begin{abstract}
This paper presents the design, implementation, and evaluation of a low-cost, low-power LoRa wireless sensor network for monitoring soil and environmental conditions beneath Icelandic turf roofs. Sensor nodes built around Arduino Pro Mini microcontrollers measure soil moisture, soil temperature, air temperature, relative humidity, and battery voltage, store records on an SD card, and communicate over 433~MHz LoRa with mesh relay functionality. A base station uploads the measurements to the ThingSpeak cloud platform for remote access. Under the defined operating cycle, a node draws an average current of approximately 3.08~mA, which gives an idealized lifetime of about 142 days from a nominal capacity of 10,500~mAh. Field tests in Vellir, Hafnarfj\"{o}r\dh ur observed single-hop links at several locations within about 3~km, with RSSI values from $-103$ to $-113$~dBm. A second-order calibration of the capacitive soil moisture sensor over the 0--40\% gravimetric range gave an in-sample RMSE of approximately 3.01 percentage points, and cloud uploads took approximately 1--2~s. Four nodes deployed beneath turf roofs at Keldur recorded valid data for approximately 51 to 130 days. The results demonstrate the feasibility of low-cost monitoring of turf roofs, while highlighting the need to consider environmental conditions and long-term hardware reliability.
\end{abstract}

\begin{IEEEkeywords}
Environmental monitoring, LoRa, low-power systems, mesh networks, soil moisture, ThingSpeak, turf roofs, wireless sensor networks.
\end{IEEEkeywords}
